% P15-U1: Durrett 2.1.15--2.1.20, with complete density and normalization prerequisites.
\chapter{Convolution, Gamma Laws and Normal Laws}
\label{sec:p15-convolution-gamma-normal}

The distribution of a sum of independent variables is determined by
their two marginal laws. When densities exist, the resulting convolution
can often be calculated explicitly. This chapter gives the measure
definition first, proves the density formulas, and performs the gamma
and normal calculations in Durrett's author Version 5 of 11 January 2019
\cite[results~2.1.15--2.1.20, pp.~49--52]{DurrettPTEAuthor2019}.
Both gamma proof routes in that presentation are supplied. The required
normalization and exponential/normal moment calculations are included
\cite[examples~1.6.10--1.6.11, pp.~33--34]{DurrettPTEAuthor2019}.

The proofs use the constructed Borel Lebesgue measures, one-variable
integration, Tonelli and the independent product laws of the preceding
chapters. The one-variable affine substitution used below is proved
from interval lengths. Gaussian normalization uses that substitution
and Tonelli; no multidimensional change-of-variables or characteristic-
function theorem is presumed.

\section{Densities and one-variable affine substitution}
\label{sec:p15-convolution-density-foundations}

\begin{hbdefinition}{def:p15-probability-density}{A density relative to Borel Lebesgue measure}
For \(d\geq1\), a Borel function \(f:\mathbb R^d\to[0,\infty]\)
is a \emph{probability density} for a Borel probability measure \(\eta\)
when
\[
 \eta(A)=\int_A f\,d\lambda^d
                 \qquad(A\in\mathcal B(\mathbb R^d)).
\]
Here \(\lambda^d\) is the constructed finite product Borel Lebesgue
measure; \(\lambda^1=\lambda\). A random vector has density \(f\)
when its law has that density. Densities are nonnegative and integrate
to one, but their values need not be at most one. The value \(f(x)\)
is not the point mass \(\mathbb P(X=x)\).

For \(c\in\mathbb R\), the \emph{Dirac probability measure}
\(\delta_c\) gives a Borel set measure one if it contains \(c\)
and zero otherwise. It is the law of a variable equal to \(c\) a.s.
\index{probability density}\index{Dirac measure}
\end{hbdefinition}

\begin{hbtheorem}{thm:p15-density-integrals}{proved=annex:p15.density-integrals}{Constructing a density law and integrating against it}
Every nonnegative Borel \(f\) with \(\int f\,d\lambda^d=1\)
defines a probability law by the displayed formula. For nonnegative
Borel \(h\),
\[
 \int h\,d\eta=\int hf\,d\lambda^d,
\]
with the nonnegative product convention \(0\cdot\infty=0\).
For signed integrable \(h\), apply this formula to its positive
and negative parts and subtract their finite integrals. In particular,
integrability is equivalent to \(\int|h|f\,d\lambda^d<\infty\).

Two densities of the same law agree \(\lambda^d\)-a.e., and measurable
nonnegative a.e. replacement preserves a density law. A density is finite a.e.;
its values at the infinite-value null set may be set to zero. Every
singleton has zero probability under a density law. Consequently
\(\delta_c\) has no density relative to \(\lambda\).
\end{hbtheorem}

\begin{hbtheorem}{thm:p15-affine-lebesgue}{proved=annex:p15.affine-lebesgue}{Affine scaling of measure, integrals and densities}
Let \(c,d\in\mathbb R\) with \(c\ne0\). For every Borel set
\(A\subseteq\mathbb R\),
\[
 \lambda(cA+d)=|c|\lambda(A),
 \qquad cA+d=\{cx+d:x\in A\}.
\]
For nonnegative Borel \(h\),
\[
 \int_{\mathbb R}h(cx+d)\,d\lambda(x)
                 =\frac1{|c|}\int_{\mathbb R}h(u)\,d\lambda(u).
\]
The same identity holds for integrable signed \(h\), and integrability
on one side is equivalent to absolute integrability on the other.
Restrictions to intervals follow by including the appropriate indicator.

If \(X\) has density \(f\), then \(cX+d\) has density
\[
 t\longmapsto\frac1{|c|}f\left(\frac{t-d}{c}\right).
\]
For \(c=0\), the transformed law is \(\delta_d\), and this
density formula is not used.
\end{hbtheorem}

\section{Convolution of probability measures and densities}
\label{sec:p15-convolution-laws}

\begin{hbdefinition}{def:p15-convolution-law}{The convolution law and its cumulative function}
For Borel probability measures \(\mu,\nu\) on \(\mathbb R\),
define their \emph{convolution} by
\[
 (\mu*\nu)(A)\defeq
 (\mu\times\nu)(\{(x,y):x+y\in A\})
                       \qquad(A\in\mathcal B(\mathbb R)).
\]
It is the pushforward of the constructed product measure under addition.
For their cumulative distribution functions \(F,G\), write
\(F*G\) for the cumulative function of this convolution law.
The notation \(\int h(y)\,dG(y)\) means integration with respect
to the probability measure \(\nu\) having cumulative function \(G\);
it does not assume that \(G\) has a derivative.
\index{convolution!probability laws}
\end{hbdefinition}

\begin{hbtheorem}{thm:p15-convolution-law}{proved=annex:p15.convolution-law}{The law of an independent sum}
The convolution \(\mu*\nu\) is a probability measure. If real
\(X,Y\) are independent with laws \(\mu,\nu\) and cumulative
functions \(F,G\), then \(X+Y\) has law \(\mu*\nu\) and
\[
 \begin{aligned}
 (F*G)(z)&=\int_{\mathbb R}F(z-y)\,dG(y)\\
         &=\int_{\mathbb R}G(z-x)\,dF(x)
                     \qquad(z\in\mathbb R).
 \end{aligned}
\]
For every nonnegative Borel \(h\),
\[
 \int h\,d(\mu*\nu)
           =\int\!\int h(x+y)\,d\mu(x)\,d\nu(y).
\]
Convolution is commutative and associative, and
\((\mu*\delta_c)(A)=\mu(A-c)\); in particular \(\delta_0\)
is its identity. For independent integer-valued \(X,Y\),
\[
 \mathbb P(X+Y=n)
      =\sum_{k\in\mathbb Z}\mathbb P(X=k)\mathbb P(Y=n-k).
\]
The cumulative-function identity is Durrett's Theorem 2.1.15; the
integer identity also supplies Exercise 2.1.10.
\end{hbtheorem}

\begin{hbtheorem}{thm:p15-convolution-density}{proved=annex:p15.convolution-density}{One density is enough for a convolution density}
Let independent real \(X,Y\) have respective laws \(\mu,\nu\),
and suppose \(X\) has density \(f\). Then a density of \(X+Y\) is
\[
 h(z)=\int_{\mathbb R}f(z-y)\,d\nu(y)
     =\int_{\mathbb R}f(z-y)\,dG(y),
\]
where \(G\) is the cumulative function of \(\nu\).
If \(Y\) also has density \(g\), this becomes
\[
 (f*g)(z)\defeq\int_{\mathbb R}f(z-y)g(y)\,d\lambda(y).
\]
The function is nonnegative and Borel measurable, has integral one,
and is finite a.e.; possible infinite values on a null set are allowed
by the density definition. The symbol \(*\) here combines density
functions, while the preceding definition combines measures or their
cumulative functions. The three operations are linked by the proved laws.
\end{hbtheorem}

\begin{hbtheorem}{thm:p15-factorized-joint-density}{proved=annex:p15.factorized-joint-density}{A factorized joint density determines independent marginals}
Suppose a real random vector \((X_1,\ldots,X_n)\), \(n\geq1\),
has density
\(p(x_1,\ldots,x_n)=\prod_{i=1}^n g_i(x_i)\), where the
\(g_i\) are nonnegative Borel functions. Define
\(c_i=\int g_i\,d\lambda\). Then \(0<c_i<\infty\),
\(\prod_i c_i=1\), and the variables are independent with respective
densities \(g_i/c_i\). Thus the factors need not initially be
normalized densities. This is the criterion in source Exercise 2.1.1.
\end{hbtheorem}

\begin{hbexample}{def:p15-convolution-elementary-examples}{Point masses and a triangular density}
The sum of two independent fair binary variables has masses \(1/4\),
\(1/2\), \(1/4\) at zero, one and two, respectively. The middle
mass comes from the two disjoint outcomes \((0,1)\) and \((1,0)\).
These are point masses of an atomic law.

For independent uniform \(U,V\) on \((0,1)\), the density convolution
counts the length of \((0,1)\cap(z-1,z)\), and gives
\[
 h(z)=\begin{cases}
 z,&0<z\leq1,\\
 2-z,&1<z<2,\\
 0,&\text{otherwise}.
 \end{cases}
\]
It integrates to one, while \(\mathbb P(U+V=z)=0\) at every
fixed real \(z\). A density value is not a point probability.
\end{hbexample}

\section{Gamma normalization and the two sum proofs}
\label{sec:p15-gamma-convolution}

\begin{hbdefinition}{def:p15-gamma-density}{Gamma shape and rate}
For \(\alpha>0\), define
\[
 \Gamma(\alpha)=\int_0^\infty x^{\alpha-1}e^{-x}\,d\lambda(x).
\]
The defining integral is over \((0,\infty)\), and the integrand's
value at zero is set to zero. The next theorem proves the integral
is finite and positive.
For shape
\(\alpha>0\) and rate \(\lambda>0\), define
\[
 g_{\alpha,\lambda}(x)=\begin{cases}
 \displaystyle\frac{\lambda^\alpha}{\Gamma(\alpha)}
                 x^{\alpha-1}e^{-\lambda x},&x>0,\\
 0,&x\leq0.
 \end{cases}
\]
The associated law is denoted \(\operatorname{gamma}(\alpha,\lambda)\),
in the source's shape/rate convention. In this displayed density,
\(\lambda\) is a positive real parameter, not the Lebesgue measure
appearing after an integral's \(d\lambda(x)\); the argument distinguishes
the two roles. We set the density at zero to zero, so no value of
\(0^{\alpha-1}\) is needed. Changing a singleton density value
does not change the law. Rate \(\lambda\) corresponds to scale
\(1/\lambda\).
\index{gamma distribution}\index{gamma function}\index{rate parameter}
\end{hbdefinition}

\begin{hbtheorem}{thm:p15-gamma-normalization}{proved=annex:p15.gamma-normalization}{Gamma integrals, normalization and moments}
For \(\alpha>0\), \(0<\Gamma(\alpha)<\infty\),
\(\Gamma(1)=1\), and
\(\Gamma(\alpha+1)=\alpha\Gamma(\alpha)\). Consequently
\(\Gamma(k+1)=k!\) for integers \(k\geq0\).
For every \(\alpha,\lambda>0\), \(g_{\alpha,\lambda}\) is a
probability density. A variable \(X\) with this law satisfies
\[
 \mathbb E[X^k]=\frac{\Gamma(\alpha+k)}
                       {\lambda^k\Gamma(\alpha)}\quad(k=0,1,2,\ldots),
 \qquad
 \mathbb E X=\frac\alpha\lambda,\qquad
 \operatorname{Var}(X)=\frac\alpha{\lambda^2}.
\]
The case \(\alpha=1\) is the exponential law with rate \(\lambda\),
density \(\lambda e^{-\lambda x}\) for \(x>0\), moments
\(k!/\lambda^k\), mean \(1/\lambda\) and variance \(1/\lambda^2\).
\end{hbtheorem}

\begin{hbdefinition}{def:p15-beta-integral}{The beta integral}
For \(\alpha,\beta>0\), put
\[
B(\alpha,\beta)=\int_0^1
          u^{\alpha-1}(1-u)^{\beta-1}\,d\lambda(u).
\]
This integral is over \((0,1)\); its endpoint power values are set
to zero. The same zero-endpoint convention applies to the rescaled
power integrals below.
Once its finite positive normalization is proved below, the
\(\operatorname{beta}(\alpha,\beta)\) law has density
\(u^{\alpha-1}(1-u)^{\beta-1}/B(\alpha,\beta)\) for \(0<u<1\),
and zero elsewhere. Endpoint density values are set to zero.
\index{beta integral}\index{beta distribution}
\end{hbdefinition}

\begin{hbtheorem}{thm:p15-beta-gamma-convolution}{proved=annex:p15.beta-gamma-convolution}{The beta-gamma identity and independent gamma sums}
For \(\alpha,\beta>0\), the beta integral is finite and positive,
is symmetric in its parameters, and
\[
 B(\alpha,\beta)=\frac{\Gamma(\alpha)\Gamma(\beta)}{\Gamma(\alpha+\beta)}.
\]
For \(z>0\), the source's rescaling identity is
\[
 \int_0^z(z-y)^{\alpha-1}y^{\beta-1}\,d\lambda(y)
       =z^{\alpha+\beta-1}B(\beta,\alpha).
\]
If independent \(X,Y\) have laws
\(\operatorname{gamma}(\alpha,\lambda)\) and
\(\operatorname{gamma}(\beta,\lambda)\), with the same rate
\(\lambda>0\), then \(X+Y\) has law
\(\operatorname{gamma}(\alpha+\beta,\lambda)\).
For a nonempty finite mutually independent family with this common rate,
the shapes add. In particular the sum of \(n\geq1\) independent
rate-\(\lambda\) exponential variables has law
\(\operatorname{gamma}(n,\lambda)\).

The beta density is normalized, and its mean and variance are
\[
 \frac\alpha{\alpha+\beta},\qquad
 \frac{\alpha\beta}{(\alpha+\beta)^2(\alpha+\beta+1)}.
\]
The annex gives both the source's normalization argument for the gamma
convolution and its separate nonnegative double-integral argument.
\end{hbtheorem}

\section{Normal normalization and the independent sum calculation}
\label{sec:p15-normal-convolution}

\begin{hbdefinition}{def:p15-normal-law}{Normal mean and variance, including a declared degenerate case}
For mean \(\mu\in\mathbb R\) and variance \(a>0\), the normal law
\(\operatorname{normal}(\mu,a)\) has density
\[
 \phi_{\mu,a}(x)=\frac1{\sqrt{2\pi a}}
                  \exp\left(-\frac{(x-\mu)^2}{2a}\right).
\]
The second parameter is variance, as in Durrett's presentation; its
standard deviation is \(\sqrt a\). The standard normal has
\(\mu=0,a=1\). The normalization and moment claims are proved next.

We additionally define \(\operatorname{normal}(\mu,0)=\delta_\mu\).
This is a declared extension of the source's positive-variance density
regime. It has no Lebesgue density and is never inserted into the formula
containing \(1/\sqrt a\). Negative variances are not permitted.
\index{normal distribution}\index{normal distribution!degenerate}
\end{hbdefinition}

\begin{hbtheorem}{thm:p15-normal-normalization}{proved=annex:p15.normal-normalization}{The Gaussian integral and the normal moments}
\[
 \int_{\mathbb R}e^{-x^2/2}\,d\lambda(x)=\sqrt{2\pi}.
\]
Thus \(\phi_{\mu,a}\) is a probability density for every
\(\mu\in\mathbb R,a>0\). Its absolute first and second moments
are finite, its mean is \(\mu\), and its variance is \(a\).
For a standard normal \(Z\), \(\mathbb E|Z|=\sqrt{2/\pi}\),
\(\mathbb E Z=0\), and \(\mathbb E Z^2=1\).
These mean and variance statements also hold for the declared law
\(\operatorname{normal}(\mu,0)\).
\end{hbtheorem}

\begin{hbtheorem}{thm:p15-normal-sums}{proved=annex:p15.normal-sums}{Affine normal transformations and independent sums}
If \(X\) has law \(\operatorname{normal}(\mu,a)\), \(a\geq0\),
then for real \(c,d\), \(cX+d\) has law
\(\operatorname{normal}(c\mu+d,c^2a)\).
If independent \(X,Y\) have laws
\(\operatorname{normal}(\mu,a)\) and
\(\operatorname{normal}(\nu,b)\), with \(a,b\geq0\), then
\[
 X+Y\text{ has law }\operatorname{normal}(\mu+\nu,a+b).
\]
For positive \(a,b\), the proof is the source's full square-completion
convolution calculation, with every normalization constant retained.
The zero-variance cases follow separately from the Dirac convention.
In a nonempty finite mutually independent family of normal variables, the means
and variances add.
\end{hbtheorem}

\begin{hbtheorem}{thm:p15-convolution-calculations-boundaries}{proved=annex:p15.convolution-calculations-boundaries}{Worked laws and the independence and parameter boundaries}
\par\nobreak
\begin{enumerate}
\item Two independent rate-two exponential variables have sum density
      \(4z e^{-2z}\) for \(z>0\), mean one and variance \(1/2\).
      The beta law with parameters \(2,3\) has density
      \(12u(1-u)^2\) on \((0,1)\).
\item Independent laws \(\operatorname{normal}(1,4)\) and
      \(\operatorname{normal}(-2,9)\) have sum
      \(\operatorname{normal}(-1,13)\), with density
      \((26\pi)^{-1/2}\exp(-(z+1)^2/26)\).
\item If independent exponential variables instead have rates one and
      two, their sum density is \(2(e^{-z}-e^{-2z})\) for \(z>0\).
      Its mean is \(3/2\), its variance is \(5/4\), and it is
      not \(\operatorname{gamma}(2,\lambda)\) for any rate
      \(\lambda>0\). The common-rate hypothesis matters.
\item For a standard normal \(Z\), take \(X=Z,Y=-Z\).
      Each has standard normal law, but their sum is zero a.s., with
      law \(\delta_0\), rather than \(\operatorname{normal}(0,2)\).
      Thus the normal-sum formula cannot discard independence.
\item For \(\alpha\leq0\), the same gamma integral
      diverges at zero. At rate zero the exponential expression
      \(\lambda e^{-\lambda x}\) is identically zero and cannot
      be a probability density. The declared gamma parameter domain
      excludes these cases.
\end{enumerate}
The elementary binary and uniform convolutions in
Example~\ref{def:p15-convolution-elementary-examples} are also verified
in the annex.
\end{hbtheorem}

\begin{proofboundary}{def:p15-convolution-next}{The next infinite construction}
The source's convolution, gamma and normal sum presentation is now
supplied with its measure and normalization prerequisites.
\Cref{sec:p15-countable-product-measures} next proves Kolmogorov's
extension theorem, Theorem 2.1.21 and Appendix A.3, and the
measurable-space representation in Theorem 2.1.22 with its explicit
Borel image and inverse. Other discrete convolution families, the
remaining independence exercises and later limit theorems require
their further treatments.
\end{proofboundary}
